\documentclass[11pt,a4paper]{article}
\usepackage[utf8]{inputenc}
\usepackage[T1]{fontenc}
\usepackage{amsmath,amssymb,amsfonts,amsthm}
\usepackage[margin=0.88in]{geometry}
\usepackage{booktabs,array,tabularx,longtable}
\usepackage{graphicx}
\usepackage{xcolor}
\usepackage{enumitem}
\usepackage{hyperref}
\usepackage{microtype}
\usepackage{cleveref}

\hypersetup{
  colorlinks=true,
  linkcolor=blue!45!black,
  citecolor=blue!45!black,
  urlcolor=blue!55!black
}
\setlist{nosep,leftmargin=*}
\newcolumntype{L}[1]{>{\raggedright\arraybackslash}p{#1}}
\newcolumntype{Y}{>{\raggedright\arraybackslash}X}
\newtheorem{proposition}{Proposition}
\newtheorem{assumption}{Assumption}
\newcommand{\R}{\mathbb{R}}
\newcommand{\E}{\mathbb{E}}
\newcommand{\I}{\mathbb{I}}
\newcommand{\tr}{\operatorname{tr}}
\newcommand{\wrap}{\operatorname{wrap}}
\newcommand{\diag}{\operatorname{diag}}

\title{\textbf{Methodological Research Plan}\\
Validation and Adversary-Aware Uncertainty for an Edge-Hosted Mobile Digital Twin}
\author{}
\date{}

\begin{document}
\maketitle

\begin{abstract}
This document defines a methodology for constructing, validating, and
security-testing an edge-hosted digital twin of a tracked UGV01 rover. The
twin combines encoder-derived motion, inertial telemetry, BN220 GNSS
measurements, and edge-observed communication timing in an extended Kalman
filter (EKF). Its security function scores GNSS measurements against a
pre-update prediction and studies whether residual-driven covariance
adaptation can reduce attack sensitivity. The methodology deliberately
separates three ideas that are often conflated: sensor agreement, statistical
filter consistency, and physical trajectory accuracy. Agreement and
normalized-innovation tests can be computed from the existing telemetry, but
physical accuracy requires a synchronized reference that is independent of
the sensors used by the estimator. A calibrated camera and AprilTag board are
therefore specified as an affordable evaluation instrument, not as an input to
the detector. The document derives the state estimator, uncertainty policies,
innovation detector, direction-dependent detectability boundary, and paired
covariance-feedback analysis. It then defines benign collection, ground-truth
validation, attack replay, statistical reporting, reproducibility, and
stage-gated execution. No experimental result is claimed here; the purpose is
to freeze a defensible research design before confirmatory evaluation.
\end{abstract}

\section{Research Objective and Scope}

\subsection{Research problem}

An edge-hosted digital twin receives telemetry from a physical rover and uses
that telemetry to maintain a synchronized estimate of rover position, heading,
uncertainty, packet health, and attack-detection capability. Edge execution
permits richer estimation and security analysis than the rover's embedded
controller, but it also introduces packet delay, jitter, loss, staleness, and
clock-offset uncertainty. Meanwhile, tracked locomotion introduces slip,
surface-dependent turning, encoder-scale error, and kinematic-model mismatch.

The security problem is not merely whether a large GNSS jump can be detected.
The central question is whether an uncertainty model that adapts to real
operating variation can remain statistically calibrated without allowing an
attacker-controlled GNSS residual to inflate the covariance that normalizes
the detector. This creates a feedback path:
\[
\text{attacked GNSS}\rightarrow\text{residual feature}
\rightarrow\mathbf Q_k\rightarrow\mathbf S_k
\rightarrow\text{normalized score}.
\]

\subsection{Research questions}

\begin{enumerate}
  \item How accurately does the digital twin reproduce the rover trajectory
  when evaluated against an independent, time-synchronized physical reference?
  \item Are the EKF innovations and covariance estimates statistically
  consistent across speed, surface, route phase, and communication condition?
  \item How does direction-dependent innovation covariance determine the GNSS
  manipulation required for a target detection probability?
  \item Can residual-coupled covariance adaptation be influenced by semantic
  GNSS manipulation, and can GPS-independent or evidence-gated adaptation
  remove that direct influence path?
  \item How do estimator accuracy, false-alarm behavior, detection delay, and
  tolerance-exceeding physical divergence change across benign and attacked
  conditions?
\end{enumerate}

\subsection{Claim hierarchy}

The study uses the claim hierarchy in \cref{tab:claim-hierarchy}. A lower tier
cannot be silently promoted to a higher tier.

\begin{table}[ht]
\centering
\caption{Evidence required for each type of claim.}
\label{tab:claim-hierarchy}
\small
\begin{tabularx}{\linewidth}{L{0.17\linewidth}Y Y}
\toprule
\textbf{Claim type} & \textbf{Required evidence} & \textbf{What it does not establish}\\
\midrule
Implementation correctness &
Unit tests, packet-schema checks, replay invariants, finite outputs, and timing-field preservation. &
Physical localization accuracy or security effectiveness.\\
Sensor agreement &
GPS--prediction residuals, encoder--IMU heading disagreement, stationary variation, and loop closure. &
Which sensor is correct when sensors disagree.\\
Statistical consistency &
NIS distribution, innovation coverage, temporal whiteness, and, with ground truth, NEES. &
Small physical error; a consistent estimator may be imprecise.\\
Physical accuracy &
Error against an independent calibrated trajectory reference with quantified reference uncertainty. &
Attack robustness outside the tested threat model.\\
Security effectiveness &
Frozen-policy benign validation and controlled attacks evaluated against independent physical and cyber references. &
Protection from attackers granted capabilities outside the stated trust boundary.\\
\bottomrule
\end{tabularx}
\end{table}

\section{Scientific Context and Related Work}

\subsection{Digital-twin credibility}

A digital twin is more than a visualization of telemetry: it is a digital
representation synchronized with a physical asset for observation,
prediction, diagnosis, or decision support. NIST's digital-twin credibility
work emphasizes verification, validation, and uncertainty quantification
(VVUQ) as prerequisites for trustworthy use \cite{shao2023credibility,nistDT}.
This motivates treating software verification, estimator consistency, and
physical validation as distinct activities in the present study.

\subsection{Localization accuracy and external reference systems}

Robotics papers ordinarily evaluate localization against a reference that is
more accurate than, and independent of, the estimator inputs. The TUM RGB-D
benchmark used synchronized motion-capture poses and reported trajectory-error
metrics \cite{sturm2012tum}. EuRoC provided synchronized visual--inertial data,
intrinsic and extrinsic calibration, and reference poses from a laser tracker
or motion-capture system \cite{burri2016euroc}. These benchmarks illustrate the
reason GNSS, encoders, and IMU cannot simultaneously be estimator inputs and
unqualified ``ground truth.''

AprilTag provides identifiable planar fiducials whose corner geometry supports
camera-relative pose estimation \cite{olson2011apriltag,wang2016apriltag2}.
Camera calibration follows established planar-calibration methods
\cite{zhang2000calibration}. AprilTag is not intrinsically ground truth: it
becomes a reference only after camera calibration, world-frame surveying,
rigid mounting, temporal alignment, and empirical uncertainty validation.

\subsection{Kalman-filter consistency and adaptive uncertainty}

The Kalman filter and EKF propagate both a state estimate and its covariance
\cite{kalman1960,jazwinski1970}. The normalized innovation squared (NIS) is a
standard measurement-space consistency statistic; under the correctly
specified linear-Gaussian model it follows a chi-square distribution
\cite{barshalom2001}. Innovation covariance matching and related adaptive
filters estimate noise statistics from residual histories
\cite{mehra1970,akhlaghi2017}. These methods can improve tracking under
changing noise, but they also create a security concern when the adaptation
input is controllable by an adversary.

Robust innovation gates, M-estimator updates, and sequential detectors offer
important comparators. Huber-style robust filtering limits the influence of
outlying measurements \cite{gandhi2010}; CUSUM accumulates persistent
departures and has been applied to compromised-sensor detection
\cite{murguia2016}. These baselines solve related but not identical problems:
measurement rejection limits immediate corruption, while the present work
focuses on the security of the uncertainty-normalization feedback path.

\subsection{Attacks on sensor fusion}

False-data injection and state-estimation attacks have been studied broadly in
cyber-physical systems \cite{pasqualetti2013,fawzi2014}. Localization-specific
work has shown that slowly evolving GNSS manipulation can exploit dynamic
properties of multisensor fusion \cite{shen2020}. The present threat model is
narrower than an RF spoofing claim: it injects semantically valid coordinate
changes at replay time and asks what the edge estimator and detector can infer
from the resulting telemetry.

\section{System, Instrumentation, and Trust Boundary}

\subsection{Physical and edge architecture}

The physical platform is a Waveshare UGV01 tracked rover with its embedded
controller, wheel/drive encoders, IMU, and a BN220 GNSS receiver. The active
firmware path is \texttt{ugv01\_gps\_dev}. The verified GNSS wiring is preserved:
\[
\text{BN220 white}\rightarrow\text{UGV01 RX},\qquad
\text{red}\rightarrow 5\text{ V},\qquad
\text{black}\rightarrow\text{GND}.
\]
The firmware exposes GPS-only telemetry through command \texttt{T:146} and
combined telemetry through \texttt{T:147}. The latter contains base motion,
encoder, voltage, IMU, GNSS quality and coordinates, sequence, and firmware
timing fields.

An edge logger polls \texttt{T:147}, estimates session clock offset, and stores
source/sample, send, edge-arrival, queue, estimate, and alarm timestamps.
Request status, sequence gaps, stale-packet state, queue depth, and HTTP
latency are recorded. Controlled delay and jitter are applied at the edge or
during replay and do not alter rover motion or source telemetry.

\subsection{Digital-twin roles}

The term ``digital twin'' refers here to four connected functions:
\begin{enumerate}
  \item a synchronized telemetry representation of the rover and network;
  \item an EKF that predicts tracked-rover motion and updates position from GNSS;
  \item uncertainty and packet-health models that express operating conditions;
  \item a security monitor that scores pre-update GNSS innovations and reports
  its condition-dependent detection capability.
\end{enumerate}
It is an edge-hosted monitoring and security estimator, not a high-precision
autonomous controller.

\subsection{Primary trust boundary}

The primary adversary has white-box knowledge of the estimator, covariance
policy, alarm, and telemetry format. It may alter only GNSS coordinate values
presented to the replayed estimator and may choose attack direction, magnitude,
start time, and temporal profile. It may not directly alter motor commands,
encoders, IMU, packet sequence, GPS quality metadata, edge-observed arrival
time, estimator or detector code, alarm state, raw logs, or evaluation
reference. The attacked coordinates may influence the estimator through the
ordinary GNSS update and, in residual-coupled variants, through the uncertainty
adaptation path; those are precisely the effects under study.

Sender timestamps are preserved but are not assumed trustworthy merely because
they are packet fields. The primary trusted timing evidence is generated by
the edge receiver's monotonic clock. Delivery-path attacks that delay, drop,
reorder, or replay whole packets form a separately labelled extension. An
attacker that controls the edge process, alarm, or protected logs is outside
scope.

\section{Digital-Twin Accuracy and Validation}
\label{sec:accuracy}

\subsection{Three different quantities}

\paragraph{Sensor agreement.}
Let $\hat{\mathbf p}_{k|k-1}$ be the pre-update predicted position and
$\mathbf z_k^{G}$ the local GNSS position. Their residual
\[
\mathbf r_k^{G}=\mathbf z_k^{G}-\hat{\mathbf p}_{k|k-1}
\]
measures GNSS-to-model disagreement. A post-update EKF--GNSS residual is even
less independent because the GNSS sample directly helped produce the posterior.
Neither residual is physical accuracy.

Encoder--IMU heading disagreement can similarly be measured as
\[
r_k^\theta =
\wrap\!\left(\Delta\theta_k^{\mathrm{IMU}}
-\Delta\theta_k^{\mathrm{enc}}\right).
\]
It diagnoses inconsistency, slip, bias, or scale error but does not identify
which sensor is correct.

\paragraph{Statistical consistency.}
The innovation and NIS are
\[
\boldsymbol\nu_k=\mathbf z_k-\mathbf H\hat{\mathbf x}_{k|k-1},
\qquad
\mathbf S_k=\mathbf H\mathbf P_{k|k-1}\mathbf H^\top+\mathbf R_k,
\qquad
d_k=\boldsymbol\nu_k^\top\mathbf S_k^{-1}\boldsymbol\nu_k.
\]
For a correctly specified linear-Gaussian measurement model of dimension
$m=2$, $d_k\sim\chi^2_2$. NIS coverage and innovation whiteness assess whether
the filter's uncertainty describes its observed residuals. They do not reveal
absolute error because the unknown physical state is absent.

\paragraph{Physical accuracy.}
Let the independent reference pose be
$\mathbf x_k^{GT}=[x_k^{GT},y_k^{GT},\theta_k^{GT}]^\top$. Then position and
heading error are
\[
e_{p,k}=\left\|\hat{\mathbf p}_k-\mathbf p_k^{GT}\right\|_2,
\qquad
e_{\theta,k}=\left|\wrap(\hat\theta_k-\theta_k^{GT})\right|.
\]
These errors quantify accuracy only after coordinate-frame alignment,
timestamp alignment, reference validation, and exclusion of invalid reference
frames.

\subsection{Why estimator sensors are not ground truth}

GNSS cannot validate a GNSS-updated EKF without circularity. Encoders are
affected by tracked-drive slip and scale error. IMU yaw integrates bias and is
not an independent global position reference. A commanded square describes
the intended route, not the route physically travelled. Tape-measured corners
provide useful sparse checks, but they do not provide synchronized position and
heading throughout each turn and straight segment.

The existing logs can therefore support internal consistency, repeatability,
and replay-security analysis. They cannot support a statement such as
``localization RMSE is $q$ cm'' until an independent trajectory stream is
available.

\subsection{Camera--AprilTag reference}

The affordable reference system consists of a fixed camera or phone on a
tripod, surveyed stationary tags defining the world frame, and a rigid
AprilTag board attached to the rover. A calibrated camera model maps tag
corners to a six-degree-of-freedom camera-relative pose. Known world-tag
coordinates then produce rover pose in a floor-fixed frame. For planar motion,
the retained reference is $(x^{GT},y^{GT},\theta^{GT})$.

The camera system is evaluation-only. It must never update the EKF, covariance
policy, attack detector, or alarm. This preserves independence.

\subsection{Reference calibration and acceptance}

The following procedure is completed before any run is accepted:
\begin{enumerate}
  \item Print tags without rescaling and measure the black-square dimensions.
  Mount the rover tag board rigidly and measure the board-to-rover transform.
  \item Calibrate camera intrinsics and lens distortion using a printed
  ChArUco or checkerboard target \cite{zhang2000calibration}.
  \item Fix focus, exposure, resolution, frame rate, tripod pose, and route
  bounds for the session.
  \item Survey at least four non-collinear world reference points in meters.
  Estimate the camera-to-world transform using all visible references.
  \item Record static validation poses at multiple positions and headings that
  were not used to fit the transform.
  \item Report static position RMSE and maximum error, heading RMSE and maximum
  error, reprojection error, tag-detection coverage, and the spatial pattern of
  residuals.
  \item Pre-register pass/fail thresholds before prospective collection. A
  practical initial target is position RMSE no greater than $0.025$ m, heading
  RMSE no greater than $2^\circ$, median reprojection error no greater than one
  pixel, and valid pose coverage of at least $95\%$. If these targets are not
  met, report the achieved uncertainty and repair the setup before claiming
  centimetre-scale accuracy.
\end{enumerate}

The reference covariance $\boldsymbol\Sigma_k^{GT}$ is estimated from repeated
static measurements and held-out surveyed poses. When reference error is not
negligible, estimator comparisons use the combined covariance
$\mathbf P_k+\boldsymbol\Sigma_k^{GT}$ rather than treating the camera as exact.

\subsection{Time and frame alignment}

Video and telemetry clocks are aligned using a recorded synchronization event
visible in both streams or an estimated clock offset with uncertainty. Ground
truth is interpolated only across short, valid intervals. Every aligned row
stores source time, camera time, alignment offset, interpolation flag, visible
tag count, reprojection error, and validity.

The primary coordinate frame is a surveyed local world frame. Post-hoc rigid
alignment may be reported as a diagnostic, as in trajectory benchmarks
\cite{sturm2012tum}, but the primary physical-accuracy result uses the
pre-calibrated world transform so that route-scale and heading errors are not
removed by fitting.

\subsection{Accuracy and consistency metrics}

For $N$ valid synchronized samples:
\[
\mathrm{RMSE}_{p}=
\sqrt{\frac{1}{N}\sum_{k=1}^{N}e_{p,k}^2},
\quad
\mathrm{MAE}_{p}=\frac{1}{N}\sum_{k=1}^{N}e_{p,k},
\quad
\mathrm{RMSE}_{\theta}=
\sqrt{\frac{1}{N}\sum_{k=1}^{N}e_{\theta,k}^2}.
\]
The study also reports median, 90th and 95th percentile error; percentages
within $0.05$, $0.10$, $0.25$, and $0.50$ m; segment-wise cross-track error;
corner/turn-angle error; endpoint loop-closure error; and relative pose error
over fixed time and distance intervals.

With independent truth, normalized estimation error squared is
\[
\mathrm{NEES}_k=
(\hat{\mathbf x}_k-\mathbf x_k^{GT})^\top
(\mathbf P_k+\boldsymbol\Sigma_k^{GT})^{-1}
(\hat{\mathbf x}_k-\mathbf x_k^{GT}),
\]
after handling angular wrapping and consistent state dimensions. NEES tests
state-space consistency; NIS remains the measurement-space test available even
without truth.

\section{Step-by-Step Benign Methodology}

\subsection{Phase 0: verification and protocol freeze}

\begin{enumerate}
  \item Freeze firmware commit, packet schema, logger version, geometry,
  uncertainty policies, alarm rules, route definitions, and run naming.
  \item Bench-check GNSS fix validity, satellite count, HDOP, checksum counter,
  IMU units, encoder signs, sequence monotonicity, timestamp ordering, queue
  depth, and packet-loss calculation.
  \item Verify session clock-offset calibration and controlled edge-side
  delay/jitter without changing source packets or rover commands.
  \item Hash raw logs and record software environment, command line, random
  seed, and exclusion reason for every rejected run.
\end{enumerate}

\subsection{Phase 1: stationary characterization}

Collect a stationary interval on each session and surface. Quantify GNSS-valid
fraction, satellite and HDOP distributions, update rate, checksum behavior,
local-position spread, stationary apparent speed, gyro bias, yaw drift, and
accelerometer variation. This phase identifies sensor floors and initialization
requirements. It does not establish moving localization accuracy.

\subsection{Phase 2: tracked-drive calibration}

The UGV01 is treated as a tracked differential-drive platform. Vendor geometry
is the frozen nominal model unless a separately pre-registered calibration
supersedes it. The calibrated quantities are left and right meters per encoder
count, effective track width, encoder signs, heading sign, IMU yaw alignment,
and command-to-speed mapping.

Use measured straight runs in both directions for encoder scale, repeated
clockwise and counter-clockwise turns for effective track width and heading
sign, and stationary/moving segments for IMU bias and alignment. Ground-truth
camera pose separates actual rover motion from encoder-predicted motion.

\subsection{Phase 3: controlled benign routes}

The controlled route family contains:
\begin{itemize}
  \item straight forward and reverse segments for translation scale and drift;
  \item a repeated square for four known direction changes, route-phase
  comparisons, and loop closure;
  \item a figure-eight for both turning directions and changing curvature;
  \item an irregular but fixed waypoint route for generalization beyond
  geometric primitives.
\end{itemize}

The square is not assumed to be ground truth. Its purpose is repeatability:
each side, corner, and loop has a consistent experimental role. The AprilTag
trajectory, rather than the motor command, records what the rover actually did.
Random free-form motion may be retained as a stress test, but a fixed irregular
route is preferred because it permits run-to-run comparison.

\subsection{Phase 4: design and prospective corpora}

The existing benign design corpus comprises a balanced matrix of two command
speeds, two surfaces, and five trials per condition on the repeated
$0.5$ m-square route. It is used for software development, model diagnosis,
attack-profile design, and exploratory threshold work. It is not untouched
confirmatory evidence because it was visible during method development.

After the ground-truth system and policy are frozen, collect a prospective
corpus without retuning. The target is:
\[
2\ \text{surfaces}\times2\ \text{speeds}\times
3\ \text{route families}\times10\ \text{trials}=120\ \text{runs}.
\]
If resources require a smaller study, prioritize a balanced subset and report
the resulting uncertainty honestly. For a run-level false-alarm target of
$0.05$, approximately 59 independent zero-alarm runs are needed for the
one-sided exact 95\% upper bound to fall below 0.05.

\subsection{Run validity}

A run is valid only if the frozen firmware and logger are used; the route is
completed without human contact; telemetry and video span the full run;
timestamp alignment succeeds; GNSS, encoder, IMU, and packet-health fields pass
schema checks; and ground-truth coverage meets the pre-registered threshold.
Environmental interruptions, camera movement, marker loss beyond the allowed
coverage, manual repositioning, logger failures, and wrong command profiles
invalidate the run. Invalid runs remain listed with reasons and are never
silently deleted.

\section{Mathematical Framework}
\label{sec:math}

\subsection{Frames, state, and controls}

The local world frame $\mathcal W$ has $x$ and $y$ axes fixed to the measured
test area. The rover state is
\[
\mathbf x_k=
\begin{bmatrix}p_{x,k}&p_{y,k}&\theta_k\end{bmatrix}^{\top}.
\]
Let encoder increments be $\Delta n_{L,k}$ and $\Delta n_{R,k}$, with
meters-per-count factors $\kappa_L,\kappa_R$ and signs $s_L,s_R$. Track travel is
\[
\Delta s_{L,k}=s_L\kappa_L\Delta n_{L,k},\qquad
\Delta s_{R,k}=s_R\kappa_R\Delta n_{R,k}.
\]
For update interval $\Delta t_k$ and effective track width $b$,
\[
v_k=\frac{\Delta s_{R,k}+\Delta s_{L,k}}{2\Delta t_k},
\qquad
\omega_k=\frac{\Delta s_{R,k}-\Delta s_{L,k}}{b\Delta t_k}.
\]
These are effective tracked-drive parameters; they do not imply ideal
no-slip wheels.

\subsection{Nonlinear prediction}

Using midpoint integration,
\[
\bar\theta_k=\theta_{k-1}+\frac{1}{2}\omega_k\Delta t_k,
\]
\[
\mathbf f(\mathbf x_{k-1},\mathbf u_k)=
\begin{bmatrix}
p_{x,k-1}+v_k\cos\bar\theta_k\,\Delta t_k\\
p_{y,k-1}+v_k\sin\bar\theta_k\,\Delta t_k\\
\wrap(\theta_{k-1}+\omega_k\Delta t_k)
\end{bmatrix}.
\]
The EKF prediction is
\[
\hat{\mathbf x}_{k|k-1}=\mathbf f(\hat{\mathbf x}_{k-1|k-1},\mathbf u_k),
\qquad
\mathbf P_{k|k-1}=\mathbf F_k\mathbf P_{k-1|k-1}\mathbf F_k^\top+\mathbf Q_k,
\]
with
\[
\mathbf F_k=
\begin{bmatrix}
1&0&-v_k\sin\bar\theta_k\,\Delta t_k\\
0&1& v_k\cos\bar\theta_k\,\Delta t_k\\
0&0&1
\end{bmatrix}.
\]

\subsection{GNSS local coordinates and update}

For a small test area around origin $(\phi_0,\lambda_0)$, latitude and
longitude are converted to local east/north coordinates using a local tangent
plane or a documented WGS84 approximation. The chosen conversion is validated
against surveyed displacements and held fixed for a session.

The position measurement model is
\[
\mathbf z_k=\mathbf H\mathbf x_k+\mathbf v_k,\qquad
\mathbf H=\begin{bmatrix}1&0&0\\0&1&0\end{bmatrix},\qquad
\mathbf v_k\sim\mathcal N(\mathbf0,\mathbf R_k).
\]
The pre-update innovation, covariance, and gain are
\[
\boldsymbol\nu_k=\mathbf z_k-\mathbf H\hat{\mathbf x}_{k|k-1},
\quad
\mathbf S_k=\mathbf H\mathbf P_{k|k-1}\mathbf H^\top+\mathbf R_k,
\quad
\mathbf K_k=\mathbf P_{k|k-1}\mathbf H^\top\mathbf S_k^{-1}.
\]
The posterior is
\[
\hat{\mathbf x}_{k|k}=\hat{\mathbf x}_{k|k-1}+\mathbf K_k\boldsymbol\nu_k.
\]
The implementation uses the Joseph covariance form
\[
\mathbf P_{k|k}=
(\mathbf I-\mathbf K_k\mathbf H)\mathbf P_{k|k-1}
(\mathbf I-\mathbf K_k\mathbf H)^\top
+\mathbf K_k\mathbf R_k\mathbf K_k^\top
\]
with the first product interpreted as
$(\mathbf I-\mathbf K_k\mathbf H)\mathbf P_{k|k-1}
(\mathbf I-\mathbf K_k\mathbf H)^\top$.

\subsection{Operational branch and proposed security-reference branch}

The current replay implementation computes the security score from the
operational EKF's pre-update innovation. This is preferable to scoring a
post-update residual, but it is not fully GPS-independent across time: the
previous posterior may already contain earlier GNSS corrections. In the current
code, IMU and timing measurements affect uncertainty and evidence gates; IMU
yaw does not directly correct the EKF heading state.

The confirmatory architecture therefore adds a parallel security-reference
branch. It is initialized during a benign alignment interval and then
propagated from protected encoder motion, commands, IMU-derived motion
evidence, and edge timing without GNSS state correction. GNSS is scored against
this branch before any operational update. The external camera reference
remains evaluation-only and is never used by either branch.

This distinction prevents an overclaim in the current paper and defines a
testable upgrade. If attacked GNSS corrects both the monitored state and the
reference used to score it, both may move toward the same false trajectory and
conceal common-mode error. The current and parallel-branch implementations
must therefore be evaluated as separate variants.

\subsection{Implemented uncertainty policies}

Let $r_k$ be the rolling admitted GNSS/dead-reckoning residual,
$\sigma_{a,k}$ the rolling vertical-acceleration standard deviation,
$\sigma_{\omega,k}$ the rolling yaw-rate standard deviation,
$\sigma^2_{v,k}$ the rolling velocity variance, and
$\delta t_k^{mis}=|\Delta t_k^{arrival}-\Delta t_k|$.
The GPS-independent planar process standard-deviation rate is
\[
c_k=q_0+g_a\sigma_{a,k}+g_\omega\sigma_{\omega,k}
+g_v\sigma^2_{v,k}+g_t\delta t_k^{mis}.
\]
The naive residual-coupled policy is
\[
q_{xy,k}^{N}=c_k+g_r r_k,
\]
whereas the GPS-independent policy is
\[
q_{xy,k}^{I}=c_k.
\]
The heading rate is
\[
q_{\theta,k}=q_{\theta,0}
+g_{\theta\omega}\sigma_{\omega,k}
+g_{\theta v}\sigma^2_{v,k}.
\]
Thus
\[
\mathbf Q_k=
\diag\!\left[
(q_{xy,k}\Delta t_k)^2,\,
(q_{xy,k}\Delta t_k)^2,\,
(q_{\theta,k}\Delta t_k)^2
\right].
\]

Measurement uncertainty uses receiver quality and timing:
\[
\sigma_{G,k}=g_h\max(h_k,h_{\min})
\left[1+\rho_s\max(0,N_{\mathrm{nom}}-N_k)\right]
\left[1+\max(0,\Delta t_k^{arrival}-\Delta t_{\mathrm{nom}})\right],
\]
where the bracketed factors multiply the HDOP term in implementation. Then
\[
\mathbf R_k=\diag(\sigma_{G,k}^2,\sigma_{G,k}^2).
\]
All coefficients and units are versioned in
\texttt{DigitalTwin/configs/uncertainty\_policies.json}.

\subsection{Evidence-gated residual feedback}

Define independent disturbance evidence
\[
E_k=\I\!\left[
|a_{z,k}-g|\ge\tau_a
\ \lor\ |\omega_{z,k}|\ge\tau_\omega
\ \lor\ |\Delta t_k^{arrival}-\Delta t_k|\ge\tau_t
\right].
\]
Residual feedback is admitted only when
\[
G_k=E_k\,
\I[d_{k-1}\le\gamma]\,
\I[\text{packet }k\text{ is not stale}]=1.
\]
The evidence-gated planar rate is therefore
\[
q_{xy,k}^{EG}=c_k+G_k g_r r_k.
\]
The protected evidence is independent of GNSS coordinates under the primary
threat model. It is not protected from a stronger attacker that can forge IMU,
timing, or packet-health fields; that case must be labelled separately.

\subsection{Innovation detector and persistent alarm}

The pre-update NIS is
\[
d_k=\boldsymbol\nu_k^\top\mathbf S_k^{-1}\boldsymbol\nu_k.
\]
For a per-update theoretical false-alarm probability $\alpha$ under the ideal
model, $\gamma_\alpha=F_{\chi^2_m}^{-1}(1-\alpha)$. The operational alarm is a
persistent rule:
\[
\mathcal A_k=
\I\!\left[
\sum_{j=k-W+1}^{k}\I(d_j>\gamma)\ge q
\right],
\]
with $W=5$ and $q=3$ for the primary NIS-family detectors. Because overlapping
windows and temporally correlated innovations violate independent Bernoulli
assumptions, run-level thresholds and false-alarm probability are validated
empirically on complete runs rather than inferred directly from $\alpha$.

\subsection{Scoped detectability derivation}

Consider one paired, pre-update, locally linearized measurement step. Let the
clean innovation satisfy
$\boldsymbol\nu_k^0\sim\mathcal N(\mathbf0,\mathbf S_k)$, and let a semantic
GNSS attack add $\mathbf a_k=\epsilon\mathbf u$ with
$\|\mathbf u\|_2=1$. If the paired prediction and covariance are held fixed at
that step,
\[
\boldsymbol\nu_k^a=\boldsymbol\nu_k^0+\epsilon\mathbf u,
\]
and the attacked NIS follows a noncentral chi-square distribution with
\[
\lambda(\epsilon,\mathbf u)
=\epsilon^2\mathbf u^\top\mathbf S_k^{-1}\mathbf u.
\]
Let $\lambda^\star$ solve
\[
\Pr\!\left[
\chi'^2_m(\lambda^\star)>\gamma
\right]=\beta
\]
for target detection probability $\beta$. Then:

\begin{proposition}[Conditional directional detectability]
\label{prop:detectability}
Under the local paired assumptions above, the attack magnitude giving target
single-update detection probability $\beta$ in direction $\mathbf u$ is
\[
\epsilon_\beta(\mathbf u)
=\sqrt{\frac{\lambda^\star}
{\mathbf u^\top\mathbf S_k^{-1}\mathbf u}}.
\]
The largest such magnitude over unit directions is
\[
\epsilon_{\beta,\mathrm{worst}}
=\sqrt{\lambda^\star\lambda_{\max}(\mathbf S_k)}.
\]
\end{proposition}

\begin{proof}
Substitution into the noncentrality expression gives the directional result.
The maximum magnitude corresponds to the minimum Rayleigh quotient of
$\mathbf S_k^{-1}$, equal to
$1/\lambda_{\max}(\mathbf S_k)$.
\end{proof}

This is not a global guarantee for the nonlinear EKF, a multi-step adaptive
attack, or the three-of-five run alarm. It is a local sensitivity statement.
Multi-step detection probability is estimated by paired replay and prospective
experiments.

A related noise-free threshold contour is
\[
\epsilon_\gamma(\mathbf u)
=\sqrt{\frac{\gamma}{\mathbf u^\top\mathbf S_k^{-1}\mathbf u}}.
\]
It describes the deterministic NIS boundary for a pure injected offset. It is
not a guaranteed ``maximum undetectable attack'' in the presence of random
innovation, persistence, feedback, or state evolution.

\subsection{Covariance-feedback mechanism}

For paired attacked and reference replays, let
$(\boldsymbol\nu_a,\mathbf S_a)$ and
$(\boldsymbol\nu_0,\mathbf S_0)$ denote the innovation/covariance pairs. Define
\[
d_a=\boldsymbol\nu_a^\top\mathbf S_a^{-1}\boldsymbol\nu_a,\quad
d_0=\boldsymbol\nu_0^\top\mathbf S_0^{-1}\boldsymbol\nu_0,\quad
d_{a|0}=\boldsymbol\nu_a^\top\mathbf S_0^{-1}\boldsymbol\nu_a.
\]
The exact score decomposition is
\[
d_a-d_0=
\underbrace{(d_{a|0}-d_0)}_{\text{innovation change in reference metric}}
-
\underbrace{(d_{a|0}-d_a)}_{\text{normalization credit}}.
\]
The normalization credit is positive when the attacked covariance reduces the
score assigned to the attacked innovation relative to the reference metric.
NIS suppression occurs when that credit exceeds the innovation growth term.

For the implemented residual-coupled planar process variance, let
$c_k\ge0$ be the independent rate, $g_r\ge0$ the residual gain, and
$r_a,r_0\ge0$ the attacked and reference rolling residuals. The exact paired
planar variance change is
\[
\Delta Q_{xy,k}
=\Delta t_k^2 g_r(r_a-r_0)
\left[2c_k+g_r(r_a+r_0)\right].
\]
This identity proves covariance inflation only when $r_a>r_0$ under the stated
mapping. It does not by itself prove larger physical error, lower run-level
detection, or attacker success.

\subsection{Safe-adaptation requirement}

For the coordinate-only threat model, a sufficient structural condition for
removing direct attack-to-$Q$ feedback is
\[
\frac{\partial\mathbf Q_k}{\partial\mathbf z_j^{G}}=\mathbf0
\quad\text{for all attacked coordinate samples }j\le k.
\]
The GPS-independent policy satisfies this condition by construction. The
evidence-gated policy satisfies it only on updates with $G_k=0$; when $G_k=1$,
residual feedback is intentionally admitted and must be evaluated. This
condition prevents direct covariance poisoning through GNSS coordinates, but
does not guarantee accurate covariance prediction or security against
multi-sensor compromise.

\subsection{Capability reporting}

The implemented ``safe/warning/blind'' visualization is a deterministic
engineering indicator formed from score margin and detectability scale. It is
not a calibrated posterior probability and must not be described as one.
Publication reporting uses the underlying NIS, $\mathbf S_k$, directional
detectability, empirical detection probability, and independently measured
trajectory error.

\section{Threat Models, Attacks, and Baselines}

\subsection{Attack families}

\begin{table}[ht]
\centering
\caption{Attack families and evaluation status.}
\small
\begin{tabularx}{\linewidth}{L{0.15\linewidth}Y Y}
\toprule
\textbf{Family} & \textbf{Attacker action} & \textbf{Scientific role}\\
\midrule
Step bias &
Add a fixed along- or cross-track GNSS displacement after attack start. &
Estimate direction-dependent detection curves and compare with
\cref{prop:detectability}.\\
Slow drift &
Increase GNSS offset gradually at a fixed rate or under a fixed budget. &
Test persistence and residual-to-covariance feedback.\\
Coordinate freeze &
Hold the last coordinate while protected telemetry continues. &
Test temporal inconsistency and dead-reckoning divergence.\\
Value replay &
Substitute coordinates from an earlier segment. &
Test temporally plausible but stale semantic content.\\
Strategic drift &
Inject during high independently predicted uncertainty using a frozen policy. &
Measure advantage from timing an attack rather than changing its budget.\\
Delivery stress &
Delay, jitter, drop, stale, or reorder packets without modifying payload. &
Secondary robustness condition, not part of the primary coordinate-only claim.\\
Combined &
Apply a content attack under delivery stress. &
Secondary composite threat with an explicitly enlarged trust boundary.\\
\bottomrule
\end{tabularx}
\end{table}

\subsection{Comparator suite}

Comparisons use the same physical run, attack, start time, and transport
condition wherever possible:
\begin{enumerate}
  \item consecutive-GNSS jump threshold;
  \item raw GNSS-to-dead-reckoning residual;
  \item fixed-covariance NIS;
  \item robust innovation gate;
  \item Huber-weighted EKF;
  \item CUSUM of whitened innovation magnitude;
  \item standard innovation-matching adaptive covariance;
  \item naive residual-coupled adaptive covariance;
  \item frozen-clean covariance schedule as a replay oracle;
  \item GPS-independent adaptive covariance;
  \item evidence-gated residual adaptation.
\end{enumerate}
The frozen-clean schedule is not deployable; it isolates covariance feedback by
reusing the clean branch's covariance under the paired attacked measurements.

\section{Evaluation and Statistical Analysis Plan}

\subsection{Benign estimator evaluation}

Report all accuracy metrics in \cref{sec:accuracy} by run and by
surface--speed--route condition. Confidence intervals are clustered or
bootstrapped by physical run, not by treating every telemetry row as
independent. Compare fixed, GPS-independent, evidence-gated, and operational
variants on physical RMSE, heading error, cross-track error, NIS coverage,
innovation whiteness, NEES, and runtime.

\subsection{Alarm locking and false-alarm validation}

All initialization, motion gates, thresholds, persistence rules, covariance
coefficients, and exclusion criteria are frozen before the prospective corpus
is inspected. The design corpus may select policies but cannot confirm their
false-alarm rate. Prospective false alarms are reported per complete run with
an exact binomial interval and per update as a secondary correlated measure.
No anomalous benign run is excluded unless it violates a pre-registered
validity rule independent of its detector score.

\subsection{Attack campaign}

Each accepted benign run is replayed clean and attacked under identical
initialization. Attacks begin at multiple pre-specified fractions of the
post-motion horizon and are applied in both along- and cross-track directions.
Step magnitude and drift-rate grids are fixed before confirmatory analysis.
Raw logs remain immutable; the injector records every modified field.

Primary outcomes are run detection probability, first-alarm delay, false-alarm
matched detection difference, maximum pre-alarm physical position error,
duration above mission tolerances, and directional
$\epsilon_{50},\epsilon_{90},\epsilon_{95}$. Until physical ground truth is
available, the corresponding replay metric is named
\emph{tolerance-exceeding paired digital-twin divergence before alarm}, not
physical harm.

\subsection{Covariance-feedback evaluation}

For matched slow-drift replays, compare attacked/reference $\mathbf Q_k$,
$\mathbf R_k$, $\mathbf S_k$, NIS, gate activation, state divergence, and
physical error. Report paired differences with physical-run-clustered
bootstrap intervals. Three conclusions are tested separately:
\begin{enumerate}
  \item mechanism: attack increases residual-driven covariance;
  \item score effect: normalization credit suppresses NIS;
  \item operational effect: suppression changes alarms or physical tolerance
  exceedance under matched false-alarm constraints.
\end{enumerate}
Evidence for one level is not treated as evidence for the next.

\subsection{Threshold and sensitivity analysis}

Sweep detector thresholds on stored score sequences to produce
false-alarm-versus-detection curves. Repeat key summaries across multiple
physical tolerances, attack directions, start times, surfaces, speeds, and
reference-quality filters. Report censored epsilon estimates when the target
detection probability lies above the largest tested magnitude.

\subsection{Multiplicity and uncertainty}

The unit of physical independence is the rover run. Use run-clustered bootstrap
confidence intervals, paired tests for matched branches, and condition-stratified
summaries. Clearly distinguish unique attack--run--start combinations from the
larger number of detector--run evaluations. Exploratory comparisons are marked
as such; confirmatory hypotheses and primary outcomes are pre-registered.

\section{Reproducibility and Data Governance}

Each run has a structured filename encoding speed, surface, network condition,
route, attack, and trial. Raw telemetry and raw video are immutable. Generated
files include a benign manifest with SHA-256 hashes, aligned ground-truth CSV,
quality-control report, replay configuration, policy versions, attack manifest,
software commit, environment lock, and exact command.

Firmware wiring and flashing details remain in operator documentation rather
than the main scientific argument. The paper-facing artifact exposes one
top-level regeneration command and separate fast validation commands. Random
seeds, bootstrap count, interpolation rules, units, coordinate frames, and
invalid-run reasons are machine-readable.

\section{Methodological Research Roadmap}
\label{sec:roadmap}

\subsection{Stage 1: manuscript and claim freeze}

\textbf{Work.} Adopt the accuracy/consistency distinction, scoped mathematical
claims, threat boundary, route definitions, metrics, and claims-versus-evidence
table in this document. Remove all numerical outcomes from the methodology
manuscript.

\textbf{Exit criterion.} Every planned claim names the evidence and dataset
required to support it; no internal residual is labelled physical accuracy.

\subsection{Stage 2: ground-truth software preparation}

\textbf{Work.} Implement camera calibration, AprilTag-board detection,
camera-to-world estimation, timestamp alignment, quality flags, trajectory
export, and synthetic/prerecorded-video tests. Define the ground-truth CSV and
visual quality-control plots before collecting confirmatory runs. In parallel,
implement the proposed GPS-independent security-state branch and tests that
prove GNSS updates cannot modify its post-alignment state.

\textbf{Exit criterion.} The pipeline reproduces known static test poses,
rejects poor detections, preserves time provenance, and runs without accessing
estimator outputs. The security branch passes causal field-influence tests and
is versioned separately from the current operational-EKF pre-update detector.

\subsection{Stage 3: physical reference validation}

\textbf{Work.} Install the fixed camera, surveyed world tags, and rigid rover
board. Calibrate intrinsics/extrinsics, test held-out positions and headings,
quantify reference uncertainty, and freeze the setup.

\textbf{Exit criterion.} Pre-registered static position, heading, reprojection,
and coverage criteria pass. Otherwise the measured limitation is recorded and
the hardware geometry is revised before accepted collection.

\subsection{Stage 4: calibration and estimator validation}

\textbf{Work.} Collect reference-tracked straight, turn, square, and
figure-eight pilot runs. Estimate effective tracked-drive parameters, quantify
slip and IMU bias, validate GNSS local coordinates, and choose $\mathbf Q$ and
$\mathbf R$ using benign-only criteria. Compare the operational-EKF pre-update
detector with the independent security-state branch. Re-freeze the estimator
and selected branch once.

\textbf{Exit criterion.} Physical accuracy, NIS, NEES, and innovation whiteness
are reported on held-out pilot runs; any estimator revision is versioned before
prospective evaluation.

\subsection{Stage 5: prospective benign validation}

\textbf{Work.} Collect the balanced untouched corpus under the frozen
estimator and alarm. Include square, figure-eight, and fixed irregular routes
where practical. Preserve all failures and exclusion reasons.

\textbf{Exit criterion.} Physical trajectory accuracy and run-level
false-alarm probability are reported with confidence intervals, without
post-hoc threshold changes.

\subsection{Stage 6: confirmatory security campaign}

\textbf{Work.} Replay the pre-registered attacks and comparator suite against
the prospective traces. Evaluate direction, magnitude, rate, injection time,
surface, speed, and transport transfer. Run the paired covariance-feedback
decomposition.

\textbf{Exit criterion.} Detection, delay, physical tolerance exceedance,
directional epsilon, and covariance-feedback claims have run-clustered
intervals and matched controls.

\subsection{Stage 7: live edge and transfer checks}

\textbf{Work.} Compare replay and live edge execution using the same frozen
packet stream and policy. Measure compute time, queueing, packet health, and
alarm equivalence. Keep delivery-path attacks separate from the coordinate-only
primary claim.

\textbf{Exit criterion.} Live and replay paths agree within documented timing
and numerical tolerances, and the edge computation meets the operational packet
rate.

\subsection{Stage 8: artifact and paper freeze}

\textbf{Work.} Freeze raw-data hashes, processed manifests, code commit,
dependencies, figures, tables, and a claims-versus-evidence matrix. Separate
methodology, confirmatory results, exploratory extensions, and limitations.

\textbf{Exit criterion.} A clean checkout can validate schemas and regenerate
paper artifacts using the documented commands.

\section{Limitations and Non-Claims}

\begin{itemize}
  \item Semantic coordinate injection reproduces detector-visible GNSS
  corruption but is not an over-the-air RF spoofing experiment.
  \item The primary threat model does not protect against compromise of the
  edge host, IMU, encoders, commands, alarm path, or reference system.
  \item A low-cost camera reference has finite accuracy and occlusion risk; its
  measured uncertainty must accompany estimator error.
  \item The EKF is a planar tracked-rover model. Its local linearization and
  no-slip approximation can fail during aggressive turns or surface slip.
  \item The detectability proposition is conditional and single-update. It
  does not replace empirical multi-step evaluation.
  \item A GPS-independent covariance policy removes one direct attack path but
  does not guarantee that its uncertainty estimate is calibrated.
  \item Results from the design corpus are developmental. Generalization and
  false-alarm claims require the prospective frozen-policy corpus.
\end{itemize}

\section{Planned Contributions}

Subject to completion of the roadmap, the work is designed to contribute:
\begin{enumerate}
  \item a VVUQ-oriented methodology that separates sensor agreement,
  statistical consistency, and physical accuracy for a mobile edge twin;
  \item a synchronized UGV01 telemetry and independent trajectory-reference
  dataset spanning motion, surface, and edge timing conditions;
  \item a scoped directional detectability analysis for innovation-based GNSS
  monitoring;
  \item a causal decomposition of attack-driven innovation growth and
  covariance normalization;
  \item GPS-coordinate-independent and evidence-gated process-uncertainty
  policies, plus a separately evaluated GPS-independent security-state branch,
  with explicit trust assumptions and matched baselines;
  \item a reproducible attack and validation protocol that distinguishes
  developmental replay evidence from prospective physical validation.
\end{enumerate}

\begin{thebibliography}{99}

\bibitem{shao2023credibility}
G. Shao, J. Hightower, and W. Schindel,
``Credibility consideration for digital twins in manufacturing,''
\emph{Manufacturing Letters}, vol. 35, pp. 24--28, 2023.
\href{https://doi.org/10.1016/j.mfglet.2022.11.009}{doi:10.1016/j.mfglet.2022.11.009}.

\bibitem{nistDT}
National Institute of Standards and Technology,
``Digital Twins for Advanced Manufacturing,'' 2026.
\url{https://www.nist.gov/programs-projects/digital-twins-advanced-manufacturing}.

\bibitem{sturm2012tum}
J. Sturm, N. Engelhard, F. Endres, W. Burgard, and D. Cremers,
``A benchmark for the evaluation of RGB-D SLAM systems,''
in \emph{Proc. IEEE/RSJ IROS}, 2012, pp. 573--580.
\href{https://doi.org/10.1109/IROS.2012.6385773}{doi:10.1109/IROS.2012.6385773}.

\bibitem{burri2016euroc}
M. Burri, J. Nikolic, P. Gohl, T. Schneider, J. Rehder, S. Omari,
M. W. Achtelik, and R. Siegwart,
``The EuRoC micro aerial vehicle datasets,''
\emph{International Journal of Robotics Research}, vol. 35, no. 10,
pp. 1157--1163, 2016.
\href{https://doi.org/10.1177/0278364915620033}{doi:10.1177/0278364915620033}.

\bibitem{olson2011apriltag}
E. Olson, ``AprilTag: A robust and flexible visual fiducial system,''
in \emph{Proc. IEEE ICRA}, 2011, pp. 3400--3407.
\url{https://april.eecs.umich.edu/papers/details.php?name=olson2011tags}.

\bibitem{wang2016apriltag2}
J. Wang and E. Olson, ``AprilTag 2: Efficient and robust fiducial detection,''
in \emph{Proc. IEEE/RSJ IROS}, 2016, pp. 4193--4198.
\href{https://doi.org/10.1109/IROS.2016.7759617}{doi:10.1109/IROS.2016.7759617}.

\bibitem{zhang2000calibration}
Z. Zhang, ``A flexible new technique for camera calibration,''
\emph{IEEE Transactions on Pattern Analysis and Machine Intelligence},
vol. 22, no. 11, pp. 1330--1334, 2000.
\href{https://doi.org/10.1109/34.888718}{doi:10.1109/34.888718}.

\bibitem{kalman1960}
R. E. Kalman, ``A new approach to linear filtering and prediction problems,''
\emph{Journal of Basic Engineering}, vol. 82, no. 1, pp. 35--45, 1960.
\href{https://doi.org/10.1115/1.3662552}{doi:10.1115/1.3662552}.

\bibitem{jazwinski1970}
A. H. Jazwinski, \emph{Stochastic Processes and Filtering Theory}.
New York, NY, USA: Academic Press, 1970.

\bibitem{barshalom2001}
Y. Bar-Shalom, X. R. Li, and T. Kirubarajan,
\emph{Estimation with Applications to Tracking and Navigation}.
New York, NY, USA: Wiley, 2001.

\bibitem{mehra1970}
R. K. Mehra, ``On the identification of variances and adaptive Kalman
filtering,'' \emph{IEEE Transactions on Automatic Control}, vol. 15, no. 2,
pp. 175--184, 1970.
\href{https://doi.org/10.1109/TAC.1970.1099422}{doi:10.1109/TAC.1970.1099422}.

\bibitem{akhlaghi2017}
S. Akhlaghi, N. Zhou, and Z. Huang,
``Adaptive adjustment of noise covariance in Kalman filter for dynamic state
estimation,'' in \emph{Proc. IEEE PES General Meeting}, 2017.
\href{https://doi.org/10.1109/PESGM.2017.8273755}{doi:10.1109/PESGM.2017.8273755}.

\bibitem{gandhi2010}
M. A. Gandhi and L. Mili,
``Robust Kalman filter based on a generalized maximum-likelihood-type
estimator,'' \emph{IEEE Transactions on Signal Processing}, vol. 58, no. 5,
pp. 2509--2520, 2010.
\href{https://doi.org/10.1109/TSP.2009.2039731}{doi:10.1109/TSP.2009.2039731}.

\bibitem{murguia2016}
C. Murguia and J. Ruths,
``CUSUM and chi-squared attack detection of compromised sensors,''
in \emph{Proc. IEEE Conference on Control Applications}, 2016, pp. 474--480.
\href{https://doi.org/10.1109/CCA.2016.7587875}{doi:10.1109/CCA.2016.7587875}.

\bibitem{pasqualetti2013}
F. Pasqualetti, F. D\"orfler, and F. Bullo,
``Attack detection and identification in cyber-physical systems,''
\emph{IEEE Transactions on Automatic Control}, vol. 58, no. 11,
pp. 2715--2729, 2013.
\href{https://doi.org/10.1109/TAC.2013.2266831}{doi:10.1109/TAC.2013.2266831}.

\bibitem{fawzi2014}
H. Fawzi, P. Tabuada, and S. Diggavi,
``Secure estimation and control for cyber-physical systems under adversarial
attacks,'' \emph{IEEE Transactions on Automatic Control}, vol. 59, no. 6,
pp. 1454--1467, 2014.
\href{https://doi.org/10.1109/TAC.2014.2303233}{doi:10.1109/TAC.2014.2303233}.

\bibitem{shen2020}
J. Shen, J. Y. Won, Z. Chen, and Q. A. Chen,
``Drift with Devil: Security of multi-sensor fusion based localization in
high-level autonomous driving under GPS spoofing,''
in \emph{Proc. 29th USENIX Security Symposium}, 2020, pp. 931--948.
\url{https://www.usenix.org/conference/usenixsecurity20/presentation/shen}.

\end{thebibliography}

\end{document}
